% Fragmento para inclusão em um capítulo de metodologia ou experimentos.
% Requer, no preâmbulo da tese: graphicx, booktabs e url (ou hyperref).
% Os fluxogramas PlantUML são considerados previamente renderizados em PNG.

\section{Experimento com a abordagem \textit{Hybrid Pipelines}}
\label{sec:experimento-hybrid-pipelines}

\subsection{Objetivo e visão geral}
\label{subsec:hybrid-visao-geral}

O experimento com a abordagem \textit{Hybrid Pipelines}%
\footnote{Código-fonte disponível em:
\url{https://github.com/kg-construction/hybrid-pipelines}.}
foi conduzido por meio de um serviço que combina um modelo de linguagem de
grande porte (\textit{large language model} --- LLM) com evidências provenientes
da Wikidata%
\footnote{Disponível em:
\url{https://www.wikidata.org/wiki/Wikidata:Main_Page}.}.
O objetivo é associar a capacidade do LLM de identificar entidades e conceitos
no texto a mecanismos de resolução de entidades e de recuperação de
informações em uma base de conhecimento estruturada.

Diferentemente do experimento baseado exclusivamente em \textit{few-shot
prompting}, a abordagem híbrida introduz uma etapa intermediária de ancoragem.
As menções extraídas do texto são associadas, quando possível, a identificadores
da Wikidata, e as relações diretas encontradas entre as entidades resolvidas são
incluídas no conteúdo enviado à etapa de construção do grafo. Além dessas
relações, o pós-processamento acrescenta conexões heurísticas entre entidades
que ocorrem no mesmo segmento textual. Desse modo, o grafo final combina a
saída do modelo, as informações recuperadas da Wikidata e regras locais da
implementação.

O processo é organizado em cinco etapas: \textit{(i)} recepção e validação da
entrada; \textit{(ii)} extração de menções; \textit{(iii)} enriquecimento com a
Wikidata; \textit{(iv)} construção e validação de um grafo RDF serializado em
Turtle; e
\textit{(v)} organização e entrega do resultado. A
Figura~\ref{fig:hybrid-processo-geral} apresenta essa visão de alto nível.

% Fonte PlantUML da Figura~\ref{fig:hybrid-processo-geral}.
% @startuml
% left to right direction
% skinparam shadowing false
% skinparam ActivityBackgroundColor White
% skinparam ActivityBorderColor Black
% skinparam ArrowColor Black
% start
% :1. Receber e validar\na entrada;
% :2. Extrair e normalizar\nmenções com LLM;
% :3. Resolver entidades\ncom Wikidata MCP;
% :4. Identificar relações\nentre QIDs resolvidos;
% :5. Montar payload com texto\ne relações; gerar RDF/Turtle com LLM;
% repeat
%   :Validar RDF com rdflib;
%   :Adicionar relações heurísticas\nde coocorrência no pós-processamento;
%   :Validar novamente com rdflib;
%   if (RDF válido?) then (sim)
%     :6. Entregar resultado\ne registrar execução;
%     stop
%   else (não)
%     if (Ainda há tentativas?\n(máximo 3)) then (sim)
%       :Reenviar ao LLM\ncom o erro de validação;
%     else (não)
%       :Retornar erro HTTP 422;
%       stop
%     endif
%   endif
% repeat while (Nova resposta RDF) is (sim)
% @enduml
\begin{figure}[htbp]
    \centering
    \caption{Visão geral do processo de criação do grafo de conhecimento.}
    \label{fig:hybrid-processo-geral}
    \includegraphics[width=0.95\linewidth]{metodo/figuras/hybrid-processo-geral.png}
    \par\smallskip
    {\footnotesize Fonte: elaborada pela autora.}
\end{figure}

\subsection{Etapas do \textit{pipeline} híbrido}
\label{subsec:hybrid-etapas}

\subsubsection{Recepção e validação da entrada}
\label{subsubsec:hybrid-entrada}

O processamento é iniciado por uma requisição HTTP do tipo \texttt{POST},
enviada à rota \texttt{/analyze}. O corpo da requisição deve conter o campo
\texttt{text}, correspondente ao conteúdo textual a partir do qual o grafo de
conhecimento será construído. A API também expõe \texttt{GET /health}, que
consulta a disponibilidade do Ollama e da Wikidata, mas essa rota não participa
da construção do grafo; ela retorna HTTP \texttt{200} somente quando os dois
componentes informam estado \texttt{ok} e HTTP \texttt{503} nos demais casos.
Opcionalmente, podem ser informados o campo
\texttt{idempotence\_key}, o número máximo de tentativas de geração do RDF
(\texttt{max\_rdf\_attempts}) e o tempo máximo de processamento
(\texttt{max\_processing\_seconds}). Apesar do
nome, a chave de idempotência é
empregada somente como identificador de correlação dos registros de execução;
ela não mantém resultados em memória nem impede o reprocessamento de uma
requisição repetida e não é incluída no objeto de resposta.

% Fonte PlantUML da Figura~\ref{fig:hybrid-receber-validar}.
% @startuml
% skinparam shadowing false
% start
% :Receber POST /analyze;
% :Interpretar o corpo JSON;
% if (text é uma string não vazia?) then (sim)
%   :Remover espaços nas extremidades;
%   :Usar idempotence_key ou gerar UUID;
%   :Calcular o prazo de execução;
%   :Registrar analyze_started quando configurado;
%   :Encaminhar texto normalizado e parâmetros;
% else (não)
%   :Retornar HTTP 400;
% endif
% stop
% @enduml
\begin{figure}[htbp]
    \centering
    \caption{Fluxo de recepção e validação da entrada.}
    \label{fig:hybrid-receber-validar}
    \includegraphics[width=0.90\linewidth]{metodo/figuras/hybrid-receber-validar.png}
    \par\smallskip
    {\footnotesize Fonte: elaborada pela autora.}
\end{figure}

Conforme a Figura~\ref{fig:hybrid-receber-validar}, o corpo JSON é inicialmente
interpretado. Para um objeto JSON, a ausência do campo \texttt{text}, o uso de
um tipo diferente de texto ou o envio de uma cadeia vazia resulta em uma
resposta HTTP \texttt{400}. O controlador não valida previamente se o valor JSON
de nível superior é um objeto; por isso, uma lista JSON pode provocar uma
exceção não tratada ao acessar o método \texttt{get}. Para entradas válidas,
removem-se os espaços nas extremidades do texto. Em seguida, o serviço define o identificador de
correlação --- gerando um UUID quando nenhum valor é fornecido --- e calcula o
prazo de execução. Esse prazo é verificado entre as etapas e seu tempo restante
é passado às chamadas ao Ollama; as consultas à Wikidata, contudo, empregam
prazo e política de repetição próprios. Assim, o campo
\texttt{max\_processing\_seconds} não interrompe uma requisição à Wikidata que
já esteja em andamento. Quando o registro de eventos está habilitado, também se
grava o início da análise. A saída da etapa é composta pelo texto normalizado e
pelos parâmetros que controlam a execução.

\subsubsection{Extração de menções}
\label{subsubsec:hybrid-extracao}

A segunda etapa identifica entidades e conceitos presentes no texto. O
procedimento sempre consulta o LLM e exige uma resposta JSON válida com ao
menos uma menção utilizável. Regras locais complementam somente padrões
lexicais específicos depois dessa validação.

% Fonte PlantUML da Figura~\ref{fig:hybrid-extrair-mencoes}.
% @startuml
% skinparam shadowing false
% start
% :Carregar prompts de sistema e extração;
% :Enviar o texto ao Ollama;
% :Interpretar a resposta JSON;
% :Validar e realinhar as menções;
% :Complementar padrões lexicais específicos;
% :Deduplicar por surface, start e end;
% :Aplicar ENTITY_MENTION_LIMIT;
% :Encaminhar menções à resolução;
% stop
% @enduml
\begin{figure}[htbp]
    \centering
    \caption{Fluxo de extração de menções de entidades e conceitos.}
    \label{fig:hybrid-extrair-mencoes}
    \includegraphics[width=0.90\linewidth]{metodo/figuras/hybrid-extrair-mencoes.png}
    \par\smallskip
    {\footnotesize Fonte: elaborada pela autora.}
\end{figure}

Como mostra a Figura~\ref{fig:hybrid-extrair-mencoes}, o serviço carrega o
\textit{prompt} de sistema e o \textit{prompt} específico da tarefa e realiza a
chamada ao Ollama. O primeiro, apresentado no Apêndice~\ref{apendiceC}, define o
comportamento geral do modelo e restringe sua resposta ao formato solicitado. O segundo, apresentado no
Apêndice~\ref{apendiceD}, recebe o texto de entrada e solicita uma estrutura JSON
estrita. Para cada menção, são requeridos a forma superficial
(\texttt{surface}), os deslocamentos inicial e final no texto (\texttt{start} e
\texttt{end}), o tipo semântico (\texttt{entity\_type}) e o grau de confiança
(\texttt{confidence}). Na interpretação da resposta, entretanto, somente uma
\texttt{surface} não vazia é obrigatória: o tipo assume \texttt{Entity} quando
ausente, e os deslocamentos e a confiança podem permanecer nulos. A estrutura
esperada é exemplificada a seguir.

\begin{verbatim}
{
  "entities": [
    {
      "surface": "...",
      "start": 0,
      "end": 5,
      "entity_type": "Entity",
      "confidence": 0.0
    }
  ]
}
\end{verbatim}

A resposta do modelo é interpretada estritamente como um objeto JSON contendo
o vetor \texttt{entities}. JSON malformado, resposta em outro formato ou ausência
de menções com forma superficial utilizável interrompem o fluxo. Falhas HTTP,
estouros de tempo ou outras exceções durante a chamada ao Ollama também são
propagadas ao controlador. Quando há menções produzidas pelo modelo, os
deslocamentos retornados são ignorados e as superfícies são buscadas
novamente, sem distinção entre maiúsculas e minúsculas, em ocorrências
sucessivas e não sobrepostas do texto normalizado. Superfícies não encontradas
são descartadas; se todas forem descartadas nesse ponto, a requisição falha.

Quando o LLM produz ao menos uma menção válida, o conjunto obtido pode ser
complementado por regras voltadas a descritores em inglês, como combinações de
modificador e termos como \textit{brand}, \textit{car}, \textit{company} e
\textit{class}, e a valores alfanuméricos posteriores a \textit{type},
\textit{class}, \textit{model}, \textit{series} ou \textit{no}. Essa
complementação integra a abordagem híbrida, mas não substitui uma extração do
LLM inválida ou vazia.

Em seguida, as menções são deduplicadas pela combinação entre forma textual,
sem distinção entre maiúsculas e minúsculas, e posições inicial e final. Essa
estratégia preserva ocorrências distintas de uma mesma expressão quando elas
aparecem em posições diferentes. Por fim, aplica-se o limite máximo configurado
em \texttt{ENTITY\_MENTION\_LIMIT}. A saída da etapa é a lista de menções
encaminhada à resolução de entidades.

\subsubsection{Enriquecimento com a Wikidata}
\label{subsubsec:hybrid-wikidata}

A terceira etapa associa as menções extraídas a entidades da Wikidata e
recupera evidências estruturadas relacionadas a esses recursos.

% Fonte PlantUML da Figura~\ref{fig:hybrid-enriquecimento-wikidata}.
% @startuml
% skinparam shadowing false
% start
% while (Há menções para resolver?) is (sim)
%   :Pesquisar a surface pelo Wikidata MCP;
%   :Manter até max(limite configurado, 10) candidatos;
%   :Selecionar por rótulo, sobreposição lexical e ordem;
%   if (Há QID e confidence é diferente de 0.5?) then (sim)
%     :Obter declarações pelo MCP;
%   endif
%   :Criar o registro WikidataEntity;
% endwhile (não)
% :Manter relações diretas entre QIDs resolvidos;
% stop
% @enduml
\begin{figure}[htbp]
    \centering
    \caption{Fluxo de enriquecimento com informações da Wikidata.}
    \label{fig:hybrid-enriquecimento-wikidata}
    \includegraphics[width=0.90\linewidth]{metodo/figuras/hybrid-enriquecimento-wikidata.png}
    \par\smallskip
    {\footnotesize Fonte: elaborada pela autora.}
\end{figure}

Conforme a Figura~\ref{fig:hybrid-enriquecimento-wikidata}, cada menção é
utilizada diretamente como texto de busca. Embora o serviço calcule o trecho da
sentença que contém a menção, esse contexto não altera a consulta enviada à
Wikidata. A busca é realizada exclusivamente por meio de um servidor baseado
no \textit{Model Context Protocol} (MCP). O cliente não implementa uma fonte
alternativa de evidência; assim, uma falha do MCP interrompe a requisição.

O candidato é selecionado pela maximização de uma tupla de pontuação. Na ordem
em que são comparados pelo código, os critérios são: igualdade exata entre a
forma superficial e o rótulo após normalização; sobreposição entre os termos da
menção e do rótulo; posição mais próxima do início da lista de candidatos; e
sobreposição entre o contexto local e o conjunto formado pelo rótulo e pela
descrição. Como a posição do candidato é comparada antes da sobreposição com o
contexto e cada posição na lista é única, a última componente da tupla não
altera a escolha entre candidatos empatados nos dois primeiros critérios.
Quando o contexto não produz termos aproveitáveis, o primeiro candidato é
selecionado diretamente. A pontuação eventualmente fornecida pela Wikidata é
armazenada no registro da entidade, mas não participa dessa escolha. Embora
\texttt{WIKIDATA\_CANDIDATE\_LIMIT} esteja
configurado como três, a implementação passa à busca o maior valor entre o
limite configurado e dez; portanto, no perfil atual, até dez candidatos podem
participar da seleção.

Depois da resolução, recuperam-se pelo MCP as declarações associadas ao
identificador selecionado. Uma falha nessa chamada é propagada sem consulta a
outra API. Qualquer menção cuja confiança seja
exatamente \texttt{0.5}, independentemente de sua origem, é resolvida, mas não
dispara a recuperação de declarações.

Mesmo quando nenhum candidato é encontrado, a menção permanece na resposta
como uma entidade não resolvida: o identificador e o IRI ficam nulos, o rótulo
recebe a forma superficial da menção e a lista de declarações permanece vazia.

Por fim, são mantidas apenas as relações diretas em que uma declaração de uma
entidade aponta para outra entidade já resolvida no mesmo texto. Cada relação
registra sujeito, propriedade e objeto, com seus identificadores e rótulos, e é
deduplicada pela combinação entre QID do sujeito, identificador da propriedade
e QID do objeto. A saída da etapa é composta pelos registros de entidades ---
resolvidas ou não ---, suas declarações e as relações diretas fundamentadas na
Wikidata.

\subsubsection{Construção e validação do grafo RDF serializado em Turtle}
\label{subsubsec:hybrid-rdf}

A quarta etapa constrói o grafo de conhecimento em RDF, serializa-o em Turtle e
verifica sua validade sintática. A construção do RDF é sempre realizada pelo LLM.

% Fonte PlantUML da Figura~\ref{fig:hybrid-criar-validar-rdf}.
% @startuml
% skinparam shadowing false
% start
% :Montar payload com texto, entidades e relações;
% repeat
%   :Solicitar RDF/Turtle ao Ollama;
%   :Limpar cercas, prefixos anteriores e notas;
%   :Validar estritamente com rdflib;
%   if (O RDF é válido?) then (sim)
%     :Aplicar pós-processamento e concluir RDF;
%     stop
%   endif
% repeat while (Ainda há tentativas?) is (sim)
% :Retornar HTTP 422;
% stop
% @enduml
\begin{figure}[htbp]
    \centering
    \caption{Fluxo de construção e validação do grafo RDF serializado em Turtle.}
    \label{fig:hybrid-criar-validar-rdf}
    \includegraphics[width=0.92\linewidth]{metodo/figuras/hybrid-criar-validar-rdf.png}
    \par\smallskip
    {\footnotesize Fonte: elaborada pela autora.}
\end{figure}

É construído um \textit{payload} JSON com o texto normalizado, a atribuição da
fonte, uma
representação compacta das entidades e as relações identificadas:

\begin{verbatim}
{
  "text": "...",
  "source_attribution": "Source: Wikidata",
  "entities": [...],
  "relationships": [...]
}
\end{verbatim}

O conteúdo substitui, em tempo de execução, o marcador
\texttt{\$\{PAYLOAD\}} do \textit{prompt} de construção do RDF, apresentado no
Apêndice~\ref{apendiceE}. A mensagem resultante é enviada ao Ollama, que deve
retornar somente o grafo em Turtle. Para cada entidade, o \textit{payload}
mantém no máximo oito declarações; as propriedades \texttt{P31}, \texttt{P279},
\texttt{P361}, \texttt{P527}, \texttt{P1889}, \texttt{P1582}, \texttt{P171} e
\texttt{P105} recebem prioridade antes desse corte.

O \textit{prompt} restringe a geração ao texto e às evidências fornecidas. Os
identificadores resolvidos devem ser preservados como recursos \texttt{wd:Q...},
e as relações recuperadas devem ser materializadas como triplas dirigidas. Para
cada recurso, exige-se um \texttt{rdfs:label} legível; no caso das entidades
resolvidas, prioriza-se a forma superficial presente no texto. Os predicados são
representados no espaço de nomes \texttt{kg:} e derivados das propriedades
disponíveis. Quando o texto expressa negação, o recurso principal também pode
receber a propriedade \texttt{kg:negated} com o valor booleano \texttt{true}.
Essa última regra é uma instrução ao LLM e não é acrescentada explicitamente
por rotinas locais.

Essas restrições descrevem o conteúdo solicitado ao modelo, mas não constituem
uma validação semântica posterior. O validador verifica se o texto pode ser
analisado como Turtle e se o grafo contém pelo menos uma tripla. O
pós-processamento remove triplas que contenham recursos \texttt{wd:Q...} fora do
conjunto de QIDs permitido, porém não elimina genericamente recursos ou triplas
no espaço de nomes \texttt{kg:} criados pelo modelo. Os rótulos ausentes são
acrescentados somente quando a implementação consegue derivar uma forma legível.
Após esse pós-processamento, o resultado é novamente validado; assim, uma
serialização que fique sem triplas ou sem relação semântica não é retornada.

Conforme a Figura~\ref{fig:hybrid-criar-validar-rdf}, a saída do modelo passa por
uma limpeza de envoltórios. Quando existe uma cerca Markdown sem rótulo ou
marcada como Turtle, TTL ou RDF, somente seu conteúdo é preservado. Também são
removidos o conteúdo anterior à primeira declaração \texttt{@prefix} e notas
iniciadas por um conjunto específico de expressões em inglês. O candidato
resultante é analisado estritamente com a biblioteca \texttt{rdflib};
declarações de prefixos sem qualquer tripla não são consideradas um grafo válido.

Quando o Turtle é válido, o pós-processamento remove triplas que contêm recursos
da Wikidata não pertencentes ao conjunto permitido, acrescenta os rótulos que
consegue derivar e materializa as relações estruturadas recuperadas. Além disso,
uma regra local
conecta pares de entidades que possuem QID, deslocamento textual e ocorrência no
mesmo segmento discursivo por relações simétricas \texttt{kg:related\_to}; quando
uma das menções é classificada como classe, conceito, objeto ou nacionalidade,
também é adicionada uma relação \texttt{kg:is} dirigida da outra entidade para
ela. Essas relações de coocorrência são inferências heurísticas da implementação
e não devem ser interpretadas como declarações recuperadas diretamente da
Wikidata.

Se o candidato não puder ser analisado e ainda houver tentativas disponíveis,
é realizada uma nova chamada à mesma etapa do modelo, cujo
\textit{prompt} inclui o erro do analisador e trechos limitados da resposta
inválida. Nenhum reparo local ou grafo substituto é tentado. O número de
tentativas do LLM é convertido para inteiro e limitado ao
intervalo de uma a três. Quando o campo é omitido, o controlador utiliza três
tentativas. Devido à expressão \texttt{max\_attempts or 3} usada no serviço, os
valores explícitos avaliados como falsos, como \texttt{0}, \texttt{false} e
\texttt{null}, resultam em três tentativas, e não em uma; valores não
convertíveis para inteiro produzem \texttt{ValueError}. O
esgotamento desse limite sem um RDF válido produz uma resposta HTTP \texttt{422},
uso específico desta implementação para sinalizar a falha de validação.

\subsubsection{Organização e entrega do resultado}
\label{subsubsec:hybrid-resultado}

A quinta etapa reúne os elementos produzidos no objeto de resposta da análise.

% Fonte PlantUML da Figura~\ref{fig:hybrid-entregar-resultado}.
% @startuml
% skinparam shadowing false
% start
% :Montar AnalyzeResponse;
% :Incluir texto normalizado;
% :Incluir entidades resolvidas e não resolvidas;
% :Incluir relações diretas;
% :Incluir RDF e Source: Wikidata;
% :Incluir resposta bruta da extração pelo LLM;
% :Registrar rdf_built quando configurado;
% :Serializar a resposta como JSON;
% :Retornar HTTP 200;
% stop
% @enduml
\begin{figure}[htbp]
    \centering
    \caption{Fluxo de organização e entrega do resultado da análise.}
    \label{fig:hybrid-entregar-resultado}
    \includegraphics[width=0.90\linewidth]{metodo/figuras/hybrid-entregar-resultado.png}
    \par\smallskip
    {\footnotesize Fonte: elaborada pela autora.}
\end{figure}

Como mostra a Figura~\ref{fig:hybrid-entregar-resultado}, a resposta contém o
texto normalizado, as entidades produzidas pela resolução --- inclusive aquelas
sem QID correspondente ---, as relações diretas identificadas, o grafo RDF
serializado em Turtle, a atribuição da fonte e a saída bruta da extração por
LLM. Como a chamada é obrigatória, esse valor sempre corresponde ao texto
devolvido pelo cliente Ollama, podendo ser uma cadeia vazia somente se a resposta
do serviço não contiver conteúdo no campo esperado. O valor de
\texttt{source\_attribution} é sempre a constante
\texttt{Source: Wikidata}, inclusive quando nenhuma menção é resolvida ou quando
o grafo contém relações heurísticas. A resposta bruta da geração do RDF não é
incluída no objeto retornado. Sua estrutura simplificada é apresentada a seguir.

\begin{verbatim}
{
  "text": "...",
  "entities": [...],
  "relationships": [...],
  "rdf": "...",
  "source_attribution": "Source: Wikidata",
  "llm": {
    "entity_extraction": "..."
  }
}
\end{verbatim}

Os eventos do processamento são registrados em JSONL quando
\texttt{ANALYZE\_LOG\_PATH} possui um valor. As chamadas ao Ollama são registradas
separadamente no arquivo CSV indicado por \texttt{OLLAMA\_CSV\_PATH}, mas apenas
depois que a geração é concluída com resposta HTTP válida. Se ocorrer falha de
escrita, o respectivo mecanismo de registro é desabilitado para as chamadas
posteriores da mesma instância. Não há um evento final específico para registrar
erros da análise.

Em uma execução bem-sucedida, a API retorna o objeto JSON com o código HTTP
\texttt{200}. Texto ausente, vazio ou de tipo inválido e valores que provoquem
\texttt{ValueError} resultam em HTTP \texttt{400}; falhas HTTP ou de conexão com
serviços externos e erros do tipo \texttt{RuntimeError}, em HTTP \texttt{502};
estouros de tempo representados por \texttt{requests.Timeout}, em HTTP
\texttt{504}; e a impossibilidade de produzir RDF válido após as tentativas
configuradas, em HTTP \texttt{422}. Exceções não contempladas pelos blocos de
tratamento do controlador podem resultar em HTTP \texttt{500}.

\subsection{Configuração experimental}
\label{subsec:hybrid-configuracao}

A Tabela~\ref{tab:hybrid-configuracao} reúne os parâmetros que interferem
diretamente na execução descrita. Os valores das variáveis de ambiente
correspondem ao arquivo \texttt{.env} atual; os demais valores são padrões
definidos pelo controlador ou pelo serviço. Como \texttt{load\_dotenv()} não
sobrescreve, por padrão, variáveis já existentes no processo, os valores
efetivos de uma execução podem diferir do arquivo quando o ambiente externo já
contém as mesmas chaves.

\begin{table}[htbp]
    \centering
    \caption{Configuração do experimento com a abordagem \textit{Hybrid Pipelines}.}
    \label{tab:hybrid-configuracao}
    \small
    \renewcommand{\arraystretch}{1.25}
    \begin{tabular}{p{4.1cm} p{6.7cm} p{2.5cm}}
        \toprule
        \textbf{Parâmetro} & \textbf{Descrição} & \textbf{Valor} \\
        \midrule
        \texttt{OLLAMA\_MODEL}
            & Modelo de linguagem utilizado nas gerações
            & \texttt{llama3.1:8b} \\
        \texttt{OLLAMA\_SEED}
            & Semente utilizada durante a geração
            & \texttt{42} \\
        \texttt{OLLAMA\_TEMPERATURE}
            & Grau de aleatoriedade da geração
            & \texttt{0} \\
        \texttt{OLLAMA\_NUM\_PREDICT}
            & Número máximo de \textit{tokens} por geração
            & \texttt{1536} \\
        \texttt{OLLAMA\_TIMEOUT\_SECONDS}
            & Prazo padrão do cliente; durante \texttt{/analyze}, prevalece o tempo global restante
            & \texttt{300} \\
        \texttt{ENTITY\_MENTION\_LIMIT}
            & Número máximo de menções enviadas à resolução
            & \texttt{16} \\
        \texttt{WIKIDATA\_CANDIDATE\_LIMIT}
            & Limite configurado; a busca aplica um mínimo interno de dez candidatos
            & \texttt{3} \\
        \texttt{WIKIDATA\_LANGUAGE}
            & Idioma enviado às consultas da Wikidata
            & \texttt{en} \\
        \texttt{WIKIDATA\_TIMEOUT\_SECONDS}
            & Prazo de cada requisição HTTP à Wikidata
            & \texttt{60} \\
        \texttt{ANALYZE\_TIMEOUT\_SECONDS}
            & Prazo global quando a requisição não informa outro valor
            & \texttt{540} (padrão) \\
        \texttt{max\_rdf\_attempts}
            & Campo da requisição; tentativas do LLM para geração do RDF
            & \texttt{1} (padrão) \\
        \bottomrule
    \end{tabular}
    \par\smallskip
    {\footnotesize Fonte: elaborada pela autora.}
\end{table}

A extração de menções sempre realiza uma chamada ao LLM. Depois do
realinhamento, da complementação e da deduplicação, até dezesseis menções são
encaminhadas sequencialmente à resolução na Wikidata. Cada chamada ao Ollama
recebe temperatura zero, semente \texttt{42} e limite de \texttt{1536} tokens de
saída. Quando a requisição não altera os padrões do controlador, o LLM recebe
três tentativas para produzir o RDF. Uma resposta inválida após esse limite
encerra a requisição.

Parâmetros opcionais, como \texttt{OLLAMA\_TOP\_K},
\texttt{OLLAMA\_TOP\_P}, \texttt{OLLAMA\_MIN\_P}, \texttt{OLLAMA\_STOP} e
\texttt{OLLAMA\_NUM\_CTX}, não foram explicitamente definidos no perfil
experimental e, por isso, não são apresentados na tabela. Também foram omitidas
configurações estritamente infraestruturais, como endereços dos serviços,
política de repetição de requisições e caminhos de arquivos de \textit{log}.

\subsection{Tratamento de falhas e reprodutibilidade}
\label{subsec:hybrid-reprodutibilidade}

A implementação envia a temperatura e a semente ao Ollama, mas não implementa
qualquer verificação de determinismo. O registro CSV do cliente contém a etapa,
o nome do modelo devolvido pelo Ollama --- ou o nome configurado, se esse campo
não estiver presente ---, o \textit{prompt}, a resposta e, quando presentes na
resposta HTTP, os campos \texttt{created\_at}, \texttt{done} e
\texttt{total\_duration}. Esse registro não inclui automaticamente o
\textit{digest} do modelo, a versão do Ollama, o \textit{commit} da aplicação ou
uma cópia integral das variáveis de ambiente.

O registro JSONL de eventos contém data e hora em UTC, chave de correlação, nome
do evento e seu conteúdo. Entre os eventos produzidos estão o início da análise,
as solicitações e respostas do LLM, os registros de entidades, as relações
encontradas, o método de validação do RDF e o RDF concluído.

A indisponibilidade do MCP interrompe a requisição e não altera a fonte de
evidência. Ainda assim, mudanças na Wikidata podem modificar candidatos e
declarações em uma execução posterior. Por essa razão, os registros produzidos
pelo serviço devem ser mantidos como parte da trilha experimental.

Em síntese, com o arquivo \texttt{.env} atual e os valores padrão da requisição,
o LLM participa da extração de menções e da geração do RDF, a Wikidata
fornece candidatos e declarações, e as rotinas locais validam e complementam o
resultado. A extração e a geração do RDF utilizam o LLM; as rotinas locais não
constroem nem corrigem o grafo. Quando o fluxo termina com HTTP \texttt{200}, o
campo \texttt{rdf} contém uma serialização Turtle produzida pelo LLM e validada
com \texttt{rdflib}. Quando todas as tentativas falham, a API retorna o erro
correspondente em vez de devolver um grafo local ou parcialmente reparado.
